% !PRINTED 533
\weberpage{533}

\begin{center}
\S.\,162.

\gesperrt{Abel'sche Gleichungen.}
\end{center}

Der Grad einer transitiven Permutationsgruppe von $m$ Ziffern
ist immer durch $m$ theilbar und also niemals kleiner als $m$.
Denn ist $P$ eine solche Gruppe und $Q_0$ der Inbegriff der Permutationen
von $P$, die die Ziffer 0 ungeändert lassen, so ist auch
$Q_0$ eine Gruppe und also ein Theiler von $P$. Nun kann man
wegen der vorausgesetzten Transitivität in $P$ ein System von
Permutationen $\pi_1, \pi_2 \ldots \pi_{m-1}$ finden, die 0 in 1, 0 in 2, \ldots
% [sic] authority prints Q\pi_1 here, omitting the subscript 0 expected from Q_0.
0 in $m-1$ überführen, und dann ist $Q\pi_1$ das System aller der
Permutationen von $P$, die 0 in 1 verwandeln. Danach ist

\weq{1}{P = Q_0 + Q_0\pi_1 + Q_0\pi_2 + \cdots + Q_0\pi_{m-1},}

also der Grad von $P$ gleich dem Producte aus $m$ und dem Grade
von $Q_0$.

\gesperrt{Die Galois'sche Gruppe einer irreducibeln Gleichung ist also niemals von niedrigerem Grade, als die Gleichung selbst.}

Eine Normalgleichung haben wir früher durch die Bestimmung
erklärt, dass sie irreducibel sei und dass jede ihrer
Wurzeln rational durch jede andere ausdrückbar sein sollte
(\S.\,145). Daraus ergiebt sich, dass die Gruppe einer Normalgleichung
sich auf die identische Gruppe reduciren muss, wenn
man eine Wurzel adjungirt; und umgekehrt ist eine Gleichung,
deren Gruppe so beschaffen ist, wenn sie zugleich irreducibel
ist, immer eine Normalgleichung.

Nun reducirt sich $P$ durch Adjunction der Wurzel $\alpha_0$ auf
$Q_0$, durch $\alpha_1$ auf $\pi_1^{-1}Q_0\pi_1$ etc. und wenn also $P$ die Gruppe
einer Normalgleichung sein soll, so ist nothwendig und hinreichend,
dass $Q_0$ die Einheitsgruppe ist, dass also durch
% [sic] authority prints an unindexed \pi before \pi_1,...,\pi_{m-1}; none is introduced here, and context calls for the identity representative.
$\pi, \pi_1, \pi_2, \ldots \pi_{m-1}$ die Gruppe erschöpft sei. Wir haben also:

\gesperrt{Damit eine irreducible Gleichung eine Normalgleichung sei, ist nothwendig und hinreichend, dass der Grad der Gruppe mit dem Grade der Gleichung übereinstimme.}

Wir betrachten hier zunächst die specielle Art von Gleichungen,
zu denen die von \gesperrt{Gauss} zuerst aufgelösten Kreistheilungs-
%CONT>>
